%hw-4.tex
%%%%%%%%%%%%%%%%%%%%%%%%%%%%%%%%%%%%%%%%%%%%%%%%%%%%%%%%%%%%%%%%%%%%%%
%%%%%%            Math 403 Homework \#4, Spring 2025            %%%%%%
%%%%%%                                                          %%%%%%
%%%%%%                 Instructor: Ezra Miller                  %%%%%%
%%%%%%%%%%%%%%%%%%%%%%%%%%%%%%%%%%%%%%%%%%%%%%%%%%%%%%%%%%%%%%%%%%%%%%
\documentclass[11pt]{article}

\oddsidemargin=17pt \evensidemargin=17pt
\headheight=9pt     \topmargin=26pt
\textheight=564pt   \textwidth=433.8pt
\parskip=1ex        %voffset=-6ex

\usepackage{amsmath,amssymb,graphicx,color,enumitem}%,setspace%,mathrsfs%,hyperref
\setlist[enumerate]{parsep=1.5ex plus 4pt,topsep=0ex plus 4pt,itemsep=0ex}
  \setlist[itemize]{parsep=1.5ex plus 4pt,topsep=-1ex plus 4pt,itemsep=-1.33ex}

\leftmargini=5.5ex
\leftmarginii=3.5ex

%for \marginpar to fit optimally
\hoffset=-1.02in
\setlength\marginparwidth{2.2in}
\setlength\marginparsep{1mm}
\newcommand\red[1]{\marginpar{\linespread{.85}\sf%
		   \vspace{-1.4ex}\footnotesize{\color{red}#1}}}
\newcommand\score[1]{\marginpar{\vspace{-2ex}\color{blue}{#1/3}}\hspace{-1ex}}
\newcommand\extra[1]{\marginpar{\color{blue}{#1/1}}\hspace{-1ex}}
\newcommand\total[2]{\marginpar{\colorbox{yellow}{\huge #1/#2}}}
\newcommand\collab[1]{\marginpar{\vspace{-11ex}\colorbox{yellow}{#1/3}}\hspace{-1ex}}
\newcommand\magenta[1]{\colorbox{magenta}{\!\!#1\!\!}}
\newcommand\yellow[1]{\colorbox{yellow}{\!\!#1\!\!}}
\newcommand\green[1]{\colorbox{green}{\!\!#1\!\!}}
\newcommand\cyan[1]{\colorbox{cyan}{\!\!#1\!\!}}
\newcommand\rmagenta[2][0ex]{\red{\vspace{#1}\magenta{\phantom{:}}\,: #2}}
\newcommand\ryellow[2][0ex]{\red{\vspace{#1}\yellow{\phantom{:}}\,: #2}}
\newcommand\rgreen[2][0ex]{\red{\vspace{#1}\green{\phantom{:}}\,: #2}}
\newcommand\rcyan[2][0ex]{\red{\vspace{#1}\cyan{\phantom{:}}\,: #2}}

%For separated lists with consecutive numbering
\newcounter{separated}

\newcommand{\excise}[1]{}
\newcommand{\comment}[1]{{$\star$\sf\textbf{#1}$\star$}}

%%%%%%%%%%%%%%%%%%%%%%%%%%%%%%%%%%%%%%%%%%%%%%%%%%%
%                                                 %
% TO ENABLE GRADING, DO NOT ALTER ABOVE THIS LINE %
%                                                 %
%%%%%%%%%%%%%%%%%%%%%%%%%%%%%%%%%%%%%%%%%%%%%%%%%%%

%new math symbols taking no arguments
\newcommand\0{\mathbf{0}}
\newcommand\CC{\mathbb{C}}
\newcommand\FF{\mathbb{F}}
\newcommand\NN{\mathbb{N}}
\newcommand\QQ{\mathbb{Q}}
\newcommand\RR{\mathbb{R}}
\newcommand\ZZ{\mathbb{Z}}
\newcommand\bb{\mathbf{b}}
\newcommand\kk{\Bbbk}
\newcommand\mm{\mathfrak{m}}
\newcommand\xx{\mathbf{x}}
\newcommand\yy{\mathbf{y}}
\newcommand\GL{\mathit{GL}}
\newcommand\FL{{\mathcal F}{\ell}_n}
\newcommand\SO{\mathit{SO}}
\newcommand\into{\hookrightarrow}
\newcommand\minus{\smallsetminus}
\newcommand\goesto{\rightsquigarrow}

%redefined math symbols taking no arguments
\newcommand\<{\langle}
\renewcommand\>{\rangle}
\renewcommand\iff{\Leftrightarrow}
\renewcommand\implies{\Rightarrow}

%to explain about LaTeX commands:
\newcommand\command[1]{\texttt{$\backslash$#1}}

%new math symbols taking arguments
\newcommand\ol[1]{{\overline{#1}}}

%redefined math symbols taking arguments
\renewcommand\mod[1]{\ (\mathrm{mod}\ #1)}

%roman font math operators
\DeclareMathOperator\im{im}

%for easy 2 x 2 matrices
\newcommand\twobytwo[1]{\left[\begin{array}{@{}cc@{}}#1\end{array}\right]}

%for easy column vectors of size 2
\newcommand\tworow[1]{\left[\begin{array}{@{}c@{}}#1\end{array}\right]}


%%%%%%%%%%%%%%%%%%%%%%%%%%%%%%%%%%%%%%%%%%%%%%%%%%%%%%%%%%%%%%%%%%%%%%
\begin{document}%%%%%%%%%%%%%%%%%%%%%%%%%%%%%%%%%%%%%%%%%%%%%%%%%%%%%%
%%%%%%%%%%%%%%%%%%%%%%%%%%%%%%%%%%%%%%%%%%%%%%%%%%%%%%%%%%%%%%%%%%%%%%


\title{\mbox{}\\[-8ex]Math 403 Homework \#4, Spring 2025\\\normalsize
	Instructor: Ezra Miller\\[-2.5ex]}
\author{Solutions by: ...your name...\\[1ex]
Collaborators: Tymon, Holly, Ben Greene\\
(1 point for each of up to 3 collaborators who also list you)}
\date{Due: 11:59pm Saturday 22 March 2025}

\maketitle

\vspace{-5ex}\collab{3}%

\noindent
\textsc{Reading assignments}\total{37}{48}%covers Lectures 16 - 22
\begin{enumerate}[itemsep=-1ex]\setcounter{enumi}{15}
\item[16--17.]\addtocounter{enumi}{2}%16. and 17.
for Tue.~4 March and Thu.~6 March
  \begin{itemize}
  \item{}%
  [Tapp, Chap.5] Lie algebras as tangent spaces to the identity
  \item{}%
  [Lax, Chap.9: Matrix-valued Functions, through \S Matrix Exponential]
  \item{}%
  [Wikipedia, \emph{Exponential map (Lie theory)}]
  \end{itemize}

\item%18.
for Tue.~18 March
  \begin{itemize}
  \item{}%
% [Stewart--Sun, \S IV.2.2] differentiability of simple eigenvalues
  [Lax, Chapter 16, p.237--240] entrywise positive matrices, Perron's theorem
  \end{itemize}

\item%19.
for Thu.~20 March
  \begin{itemize}
  \item{}%
% [Stewart--Sun, \S IV.2.2] differentiability of simple eigenvalues
  [Lax, Chapter 16, p.240--245] stochastic and nonnegative matrices,
  Frobenius thm
  \end{itemize}

\item%20.
for Tue.~25 March
  \begin{itemize}
  \item{}%
  [Treil, \S 8.5] multilinear algebra, tensor product
  \item{}%
  [Wikipedia, \emph{Tensor product}]
  \end{itemize}

\item[21--22.]\addtocounter{enumi}{2}%21. and 22.
for Thu.~27 March and Tue.~1 April
  \begin{itemize}
  \item{}%
  [Wikipedia, \emph{Exterior algebra}]
  \end{itemize}
\end{enumerate}

\noindent
\textsc{Exercises}% covers Lectures 14 - 17

\begin{enumerate}

\item\ \score{0}%1.
Show that $\Vert AB\Vert_2 \leq \Vert A\Vert_2\Vert B\Vert_2$ whenever
the product $AB$ is defined.

\textbf{Solution}

\[
\| A \|_2 = \sigma_{\max}(A)
\]

Since the spectral norm is the max streching factor of the matrix, the product
can't strecth more than the product of the individual matrices. \magenta{If a more 
rigorous proof is needed}, assume that it is false, and that there is \rmagenta{A more rigorous proof is indeed needed}
a vector $x$ such that $|ABx|> ||A||_{2}||B||_{2}x$. \red{$\|A\|\|B\|x\neq\|Ax\|\|Bx\|$}But the $|Bx|$ is bounded by
$|||B||_{2}x|$, so $ABx$ is bounded by $|||B||_{2}Ax|$,  and $Ax$ itself is bounded by 
$|||A||_{2}x|$, so we get a contradiction, because $|ABx| \leq|||B||_{2}Ax| \leq |||A||_{2}||B||_{2}x|$.

\item\ \score{2}%2.
Show that $\lim_{k\to\infty} A^k = 0$ if and only if its spectral
radius satisfies $\rho(A) < 1$. \red{Part II not quite rigorous}

\textbf{Solution}

\[
\rho(A) = \max \{ |\lambda| \mid \lambda \text{ is an eigenvalue of } A \}
\]

Claim($\rightarrow$): $\lim_{k\to\infty} A^k = 0 \rightarrow \rho(A)<1$.

Assume that $A$ has eigen value $\lambda \geq$ 1, and $x \in eigenspace(\lambda)$, then
$\lim_{k\to \infty}A^{k}x=\lambda^{k}x$, which isn't equal to 0, unless $x$ is the zero vector
which can't happen because it is an eigenvector.

Claim($\leftarrow$): $\rho(A)<1 \rightarrow \lim_{k\to\infty} A^k = 0$\\
The generalized eigen vectors of $A$ provide a basis for the entire ambient space. 
Describe every vector as a linear combination of these generalized eigenvectors. 
Since $\rho(A) < 1 \rightarrow \lim_{k\to\infty} |A^kx| = 0$, for all vectors x, since
the generalized eigen vectors all go to 0. Now take the 'Standard' basis and multiply them with 
$\lim_{k\to\infty}A^k$, they are all 0, which means the columns are 0, therefore the matrix is 0.

\item\ \score{3}%3.
Show that if $\nu$ is a consistent norm on $\CC^{n\times n}$, then
$\lim_{k\to\infty} \nu(A^k)^{\frac 1k} = \rho(A)$.

\textbf{Solution}

Let $A^{*}=\frac{A}{\rho(A)+x}$, where $x>0$. Since $\rho(A/k)=|1/k|\rho(A)$,
$\rho(A^{*})<1$. By the results of the previous question $\lim_{k\rightarrow \infty} (\nu((A^*)^{k}))<1$, since 
it is submultiplicative.

This also means also $(\nu((A^{*})^{k}))^{1/k}<1$. Decomposing to our original expression,
$(\rho(A)+x)\nu((A^{k}))^{1/k}<1$, if you take the sequence $x=1/k$, we see that this converges to $\nu(A^{k})^{1/k}=\rho(A)$ since
$\rho(A^{k})\geq \rho(A)^{k}$, for any consistent norm $\nu$.

\item\ \score{0}%4.
Prove or disprove, and salvage the statement as best you can in case
you find it's false:
%
If~$\Vert(\tilde\lambda I - A)^{-1}\Vert \geq \eta$ then there is an
eigenvalue $\lambda$ of~$A \in \CC^{n \times n}$ satisfying
$$%$$ $
  |\tilde\lambda - \lambda| \leq 2(\Vert A\Vert + \eta^{-1})\eta^{-\frac1n}.
$$%$$ $

\textbf{Solution}



\item\ \score{2}%5.
For $A \in \CC^{m \times n}$, set $\Vert A\Vert_\infty =
\displaystyle\max_{\Vert x \Vert_\infty = 1} \Vert Ax\Vert_\infty$.
Prove that $\Vert A\Vert_\infty = \displaystyle\max_i \sum_j
|a_{ij}|$. \red{Not quite}


\textbf{Solution}

Claim \#1: All the elements of $x$ are conjugates of all the elements in a row.

Assume the largest absolute value entry happens in in the row $j$.
\begin{align*}
  \|Ax\|&=\|(Ax)_{j}\|\\
  (Ax)_{j}&= A_{j}x
\end{align*}

This is a sum of values where each value is bounded by the absolute values of the 
numbers being multiplied. If you make the entries of $x$ the conjugates then you 
get the max value of this expression, and through construction we know that 
this is the max value that can be achieved since the absolute value is bounded by the 
product of the absolute values(moduleses), and since the infinity norm of $x$ is 1, then the max the 
value could be is the absolute value of the entries in each row. Which we have found a 
valid construction. Each element of $x$ is also normalized to be on the unit circle.

\item\ \score{1}%6.
Assume that the union of $m$ out of the $n$ Gerschgorin disks
$\mathcal{G}_i$ is disjoint from the other $n - m$ disks.  Prove that
the union of $m$ disks contains precisely $m$ of the eigenvalues.
Hint: how do the eigenvalues move as $A$ proceeds along a straight
line to~$\widetilde A$? \red{Not rigorous: you need an explicit deformation and an explicit use of continuity. }

\textbf{Solution}

Start off with a matrix that only contains the elements on the diagonal. The elements on the diagionals 
are the eigenvalues and the union of all the Gerschgorin disks are just the points.
Since we will take the union of $m$ of the Gerschgorin disks, find the rows from which we will get the radius.

Slowly add multiplies of the these rows subtraced the diagonal to the matrix. So we are slowly 
'moving' towards A. Since the deformation of the eigenvalues is contininous the eigenvalues can't 'jump' out 
of the union of the m disks  because if they did, they would have to be in another Gerschgorin disk, but since the 
$m$ Gerschgorin disks are disjoint from the ones, it would be left in no Gerschgorin disk, which is a contradiction.

So once you finish adding back all the elements in the rows of the Gerschgorin disk, you can start adding the other rows 
subtracted the diagonal back. By the same reasoning the eigenvalues can't leave the union of the $n-m$ disks, and in the end.
You get back A with m eigenvalues in the union of $m$ of the Gerschgorin disks and the other $n-m$ in the other disks.

Finer note about 'jumping', the union of the $m$ Gerschgorin disks is closed union of circles/spheres, and the other $n-m$ 
Gerschgorin disks are just points. Since the eigenvalues move contininously, it would have to first be on the boundary of the union of the 
$m$ Gerschgorin disks and then move towards a point(Since the other n-m Gerschgorin disks are just points). You will have a nonzero distance
between these points and the boundary since they are disjoint, so the eigenvalue would have to move this distance to get to the other points. 
But that would imply that it would be in an area where it isn't in the union of any Gerschgorin disks, which is the contradiction.

\item\ \score{1}%7.
Let $G(A)$ be the union of the Gerschgorin disks of~$A$.  Show that
the intersection $\bigcap_S G(S^{-1}AS)$ over all nonsingular
matrices~$S$ equals the spectrum of~$A$.

\textbf{Solution}

First the eigenvalues will always be in the Gerschgorin disks since change of basis doesn't change the eigen values.

Assume there is some element element except the eigenvalues in the intersection. Take the jordan canonical form of the matrix.
Since there are eigenvalues on the diagonal, the Gerschgorin disks are centered around the eigenvalues and the radius is at max 1. Since there
are only 1's in the superdiagonal. Now change the basis of so that the ones become arbigtrarily small, so $x$ can't be arbitrarily close to the \red{You need to explain how to do this}
eigenvalues, therefore only the eigenvalues are in the intersection.

\item\ \score{3}%8.
Prove that the spectrum of~$A$ is contained in $G(A) \cap G(A^\top)$.
Illustrate with the $3 \times 3$ matrix with entries $a_{ij} = i/j$.
\textbf{Solution}

The characteristic polynomial of A and $A^{T}$ are the same. So the spectrum
is in the interseciton.

$$
A=\begin{bmatrix}
  1 & 1/2 & 1/3\\
  2 & 1 & 2/3\\
  3 & 3/2 & 1
\end{bmatrix}
$$

The reduced row echelon form of the matrix \( A \) is:
\[
\text{rref}(A) =
\begin{bmatrix}
1 & 0 & -1 \\
0 & 1 & 2 \\
0 & 0 & 0
\end{bmatrix}
\]

The reduced row echelon form of the transpose of \( A \) is:
\[
\text{rref}(A^\top) =
\begin{bmatrix}
1 & 0 & -1 \\
0 & 1 & 2 \\
0 & 0 & 0
\end{bmatrix}
\]

So the eigen values are the same.

\item%\ \score{}%9.
Fix the matrix
  $$%$$ $
  A =
  \left[
  \begin{array}{@{}rrr}
     7 & -16\ &  8 \\
   -16 &  7 \ & -8 \\
     8 & -8 \ & -5 
  \end{array}
  \right].
  $$%$$ $
\begin{enumerate}%[(i)]%package {enumerate} not compatible with {enumitem}
  \item\ \score{3}%(a)
  Use Gerschgorin's Theorem to say as much as you can about the
  locations of the eigenvalues of and the spectral radius of~$A$.

  \textbf{Solution}

  The eigen values of A are within two cicles, centered around 7 of radius 24 and -5 of radius 16.
  
  \item\ \score{3}%(b)%
  Consider $D^{-1}AD$ for a diagonal matrix~$D$ to see if you can
  improve your estimates for the eigenvalue locations.

  \textbf{Solution}

  $$
  \begin{bmatrix}
    99 & 0 &0\\
    0 &9 &0\\
    0 &0 &201
  \end{bmatrix}  
  \begin{bmatrix}
  7 & -16 &8\\
  -16 &7 &-8\\
  8 &-8 &-5
  \end{bmatrix}  
  \begin{bmatrix}
  1/99 & 0 &0\\
  0 &1/9 &0\\
  0 &0 &1/201
  \end{bmatrix}
  =
\begin{bmatrix}
  7 & -176 &\frac{264}{67}\\
  \frac{-16}{11} &7 &\frac{-24}{67}\\      
  \frac{536}{33} &\frac{-536}{3} &-5
\end{bmatrix}  
  $$

  I think so? I don't get a better approximation.

  \item\ \score{3}%(c)%
  Compute the actual eigenvalues and comment on the quality of your
  estimates in~(a) and~(b).

  \textbf{Solution}

  The eigenvalues are 27, -9. -9. I guess I should have guessed 
  that the two of the eigen values are the same because the disks are the same.
  But I don't have a proof for that. I also see that on eof the eigen values is very far from the center.

  \end{enumerate}

\item\ \score{2}%10.
Find the Lie algebra $\mathfrak{so}_n$ of the special orthogonal group
$\SO_n(\RR)$.\red{This is right but it doesn't give the actual Lie algebra unless you also compute the Lie algebra of $O(n)$}

\textbf{Solution}

The Lie algebra of the orthonal group is set of the derivatives of all paths through the origin. Since any 
path along the origin inside the orthogonal group has to have determinant one(The determinant is contininous and therefore there is no
discontinious jump between the matrices of determinant 1 and -1, and the determinant of the identity is 1), the Lie algebra of the special orthogonal group is just 
the Lie algebra of the Orthogonal group, this group is also a subset of the lie algebra of the special linear group since the elements 
along the diagonal are zero in a skew matrix.

\item\ \score{3}%11.
Fix a matrix group~$G\subseteq M_n F$ over the field $F\in\{\RR,\CC\}$
and a matrix $A \in G$.  Show that the tangent space $T_A(G)$ of $G$
at~$A$ is $A\mathfrak{g}$, and show that this equals~$\mathfrak{g}A$.

\textbf{Solution}

Claim 1: A$\mathfrak{g} \subseteq T_{A}(G)$.

For every element in the Lie algera, there exists a path that crosses through the origin and has that element 
as a derivative. If you multiply this path by A, then this path crosses through A at 0, and is also in the group, the derivative 
of this new path is the old derivative times A, so A$\mathfrak{g} \subseteq T_{A}(G)$.

Claim 2: $ T_{A}(G)\subseteq A\mathfrak{g} $

Take a path in $T_{A}{G}$. At 0 the path is A. Since A is in a group it must have an inverse, multiply this path by the inverse, and then you 
get path that crosses the origin at 0, and this path is in $\mathfrak{g}$. So when you multiply it by A, you get back to the previous path, proving the lemma.

This proof also works for $\mathfrak{g}A$, so they are equal(since $Ag = T_{A}({G}) \wedge gA = T_{A}({G}) \rightarrow gA = Ag$ ).

\item\ \score{3}%12.
Prove that $\det(e^A) = e^{\mathrm{tr}(A)}$.

\textbf{Solution}

Take the jordan block of A. So $A=P^{-1}JP$.

\begin{align*}
  det(e^{A})&=det(\sum_{k=0}^{\infty}{\frac{A^{k}}{k!}})=det(P^{-1}\sum_{k=0}^{\infty}{\frac{J^{k}}{k!}}P)\\
            &det(P)*det(\sum_{k=0}^{\infty}{\frac{J^{k}}{k!}})*det(P^{-1})\\
\end{align*}

Since J, is upper diagonal $det(\sum_{k=0}^{\infty}{\frac{J^{k}}{k!}})=e^{\lambda_1 + \lambda_2 + \cdots + \lambda_n}$(J upper diagonal $\rightarrow$
the elements in the diagonal become $e^{\lambda_1},e^{\lambda_2}, \cdots $ and there are no elements in the lower diagonal), which is
just $e^{tr(A)}$

\item\ \score{2}%13.
Explain why $\mathfrak{su}_n =\mathfrak{u}_n\cap\mathfrak{sl}_n(\CC)$.
(You may use the previous exercise.)

\textbf{Solution}

Claim 1: $\mathfrak{su}_{n} \subseteq \mathfrak{u}_{n} \cap \mathfrak{sl}_{n}$

This is trivially true by defintion of the $SU$.

Claim 2: $ \mathfrak{u}_{n} \cap \mathfrak{sl}_{n} \subseteq \mathfrak{su}_{n}$

Take $A \in \mathfrak{u}_{n} \cap \mathfrak{sl}_{n}$.

Take the path realization of A. This path is in both $U_n$ and $SL_{n}$(By lecture 16 and lecture 17) and therefore 
is also in $SU_{n}$. So $e^{tA}$ is both in $U_n$ and $SL_n$.\red{You've gotta say a little more}

\item\ \score{3}%14.
Let $\gamma: ( -\varepsilon, \varepsilon ) \to \GL_n(F)$ be a
differentiable path, where $F \in \{\RR,\CC\}$.  Use Cramer's rule (or
some other method) to show that the inverse path $t \mapsto
\gamma(t)^{-1}$ is differentiable.

\textbf{Solution}

\begin{align*}
  A^{-1}=\frac{1}{det(A)}adj(A)
\end{align*}

Since the determinant is never zero and the cofactor and determinant are a rational function. The
derivative is always defined.

\end{enumerate}



%%%%%%%%%%%%%%%%%%%%%%%%%%%%%%%%%%%%%%%%%%%%%%%%%%%%%%%%%%%%%%%%%%%%%%
\end{document}%%%%%%%%%%%%%%%%%%%%%%%%%%%%%%%%%%%%%%%%%%%%%%%%%%%%%%%%
%%%%%%%%%%%%%%%%%%%%%%%%%%%%%%%%%%%%%%%%%%%%%%%%%%%%%%%%%%%%%%%%%%%%%%


%%% Various Emacs customizations:
%%% Local Variables:
%%% fill-column: 70
%%% indent-tabs-mode: t
%%% TeX-electric-sub-and-superscript: nil
%%% TeX-brace-indent-level: 0
%%% LaTeX-indent-level: 0
%%% LaTeX-item-indent: 0
%%% TeX-master: t
%%% End:
